% Template for Group T master thesis paper
%          Adapted by Koen Eneman from the ICASSP 2008 conference template (http://www.icassp2008.org/)
%					 And then adapted again by Jeroen Buijs
%          Uses GroepTcls.cls and apacite.bst
% --------------------------------------------------------------------------
\documentclass[a4paper]{article}
\usepackage{StyleEN}

\usepackage{hyperref}
\usepackage{flushend}

%%%%%%%%%%%%%%%%%%%%%%%%%%%%%%%%%%%%%%%%%%%
% Addition by Jeroen Buijs, November 5, 2008%
%%%%%%%%%%%%%%%%%%%%%%%%%%%%%%%%%%%%%%%%%%%
% To include EPS with pdflatex (needed for the logos!)
% Don't forget to include --shell-escape in the Latex->PDF output profile (Build -> Define Output Profile -> "Latex->PDF" -> "`Command line arguments to pass to the computer"')
%Notice that the two logos on top of the page are .eps, so you have .eps figures anyway.
\newif\ifpdf
\ifx\pdfoutput\undefined
    \pdffalse
\else
    \pdftrue
    \pdfoutput=1
\fi

\ifpdf
\usepackage{graphicx}
\usepackage{epstopdf}
\epstopdfsetup{suffix=}
\usepackage{epsfig}
\DeclareGraphicsRule{.eps}{pdf}{.pdf}{`epstopdf #1}
\pdfcompresslevel=9
\else
\usepackage{graphicx}
\usepackage{epsfig}
\fi

\usepackage{amsmath,amssymb}
\usepackage[english]{babel}
\usepackage[numberedbib]{apacite}

% Example of making definitions of things you can use further on
% --------------------
\def\x{{\mathbf x}}
\def\L{{\cal L}}


% Title.
% ------
\title{Modeling and Optimization of the Control Algorithm of\\ Selective Laser Sintering Machines}


% Authors
% -------
\author{Name Student One\supit{*}, Name Student Two\supit{*} \\ \vspace*{0.25cm}
            \small{\supit{*}Master student $<$Discipline$>$, GROUP T - Leuven Engineering College, Vesaliusstraat 13, 3000 Leuven \\ \vspace*{1.5cm}
            Supervisor(s): $<$Names$>$ \\            
            Unit $<$Unit Name$>$, GROUP T - Leuven Engineering College, Vesaliusstraat 13, 3000 Leuven, $<$\href{mailto:First.Last@groept.be}{First.Last@groept.be}$>$} \\ \vspace*{1em}
            Co-supervisor(s): $<$Names$>$\\
            $<$Department,$\ldots>$, $<$Company$>$, $<$Company Address$>$\\
            }


\begin{document}

\begin{center}
{
% POGING 1 - 1 rij met logo's en tekst%
%%%%%%%%%%%%%%%%%%%%%%%%%%%%%%%%%%%%%%%
\makeatletter
\vskip -140\p@
\begin{tabular}{p{0.2\textwidth}p{0.5\textwidth}p{0.2\textwidth}}
  \vspace{-2.5em} {\centering \includegraphics[width=0.2\textwidth]{GroupTEngineering_p2.eps} \par}
	& \vspace{-2.5em}
	{\small \centering 
	  Paper submitted within the scope of the Master's Thesis \\ Master of Industrial Sciences \\ GROUP T - Leuven Engineering College - 2012-2013 \par}
	& \vspace{-2.5em} {\centering \includegraphics[width=0.13\textwidth]{AssociatieLogo.eps} \par} %\\ \hline
\end{tabular} \par}
\vspace*{0.5cm} \centering{\emph{GROUP T does not guarantee error-free content of this paper.}}\\
{\vspace{5em} 
\makeatletter
\bf \huge{\@title} \par} \vskip 1.5em 
\makeatletter
{\large \@author}
\end{center}
\vspace*{5em}

\centerline{\bf ABSTRACT} \vspace*{0.2em}
\noindent Wireless Sensor Networking (WSN) is an emerging technology that is shaping the world around us. However, as WSNs proliferate, connecting these nodes becomes more and more challenging. 
In this work we designed, implemented and configured a ZigBee WSN test-bed and a middleware communication layer to extract the sensor data from the network. This paper details the challenges we faced and the solutions we developed for them. These challenges are network stability, latency, energy consumption, quality of service and both event and query-based sensor acquisition. All these challenges have been implemented and tested with Libelium's first generation Waspmote. A stable monitoring network consisting of at least five nodes, including one weather station, was deployed at Group T, campus Vesalius. This network responds to several commands sent through a web interface. Developing this network has convinced us that ZigBee is well suited for networks that share our objectives. We conclude that ZigBee is a promising protocol with many possible applications. On the other hand, developing the network also leads us to conclude that Libelium's platform still faces practical limitations when extended with ZigBee. 
%The abstract should be viewed as a miniature version of the paper. Since potential readers should be able to make their decision on the personal relevance based
%on the Abstract, the Abstract should clearly tell the reader what information he can expect
%to find in the paper. The Abstract should (1) state the principal objectives and scope of the
%investigation, (2) describe the methods employed, (3) summarize the results, and (4) state the
%principal conclusions. Most of the Abstract will be written in the past tense, because it refers
%to work done. The Abstract should never give any information or conclusion that is not stated
%in the paper. The authors should always keep in mind, that the Abstract is the most frequently
%read part of a paper. It should contain at most 200 words. Today, the Abstract is generally
%printed in one paragraph.  Define all symbols used in the Abstract. Do not cite references in
%the Abstract. Do not include or refer to tables and figures.
\\ \vspace*{2em}

\centerline{\bf Keywords} \vspace*{0.2em}
\noindent
Internet of Things, Libelium Waspmote, Wireless Sensor Network, ZigBee

%\pagebreak
\twocolumn

\section{Introduction}
\label{sec:intro}
Start the introduction at the top of the second page. Notice that the LaTeX template uses centered heading, whereas the Word template suggests left flushed headings. It's really a matter of taste, feel free to adapt.

The purpose of the Introduction should be to supply sufficient background information to the reader to understand and evaluate the results of the present study. The reader should not have to consult other publications in order to understand your Master's Thesis. 

Above all, state briefly and clearly your purpose in writing the paper. Choose references carefully to provide important background information. Much of the Introduction should be written in the present tense. 

Suggested rules for a good Introduction are as follows: (1) The Introduction should present first, with all possible clarity, the nature and scope of the problem investigated. (2) It should briefly review the pertinent literature to orient the reader. (3) It should state the method of the investigation. If deemed necessary, the reasons of the choice of a particular method should be stated. (4) It should state the principal results of the investigation. (5) It should state the principal conclusions suggested by the results. Do not keep the reader in suspense; let the reader follow the development of the evidence. 

\section{Materials and Methods}
\label{sec:matmeth}
In the first section of the paper, the Introduction, you stated the methodology. If necessary, you defended the reasons for your choice of a particular method.

\subsection{Materials}
Now, in Materials and Methods, you must give full details. Most of this section will be written
in the past tense. 

\subsubsection{Main purpose}
The main purpose of the Materials and Methods section is to describe (and if
necessary to defend) the experimental design and to provide enough detail so that a competent worker
can repeat the experiments.
\par 
\vspace*{1ex} \noindent
\textbf{Other purposes.} Other purposes include providing information that will let readers judge
the appropriateness of the experimental methods (and thus the probable validity if the findings)
and that will permit assessment of the extent to which the results can be generalized.
\par 
\vspace*{1ex} \noindent
\textbf{Just to make sure.} So notice that we don't define a fourth level of subheading! In fact, you should try to avoid deep levels like this, but if really needed, you can fabricate your own header like shown here.

\subsubsection{Research Materials}
For Materials, include the exact technical specifications and quantities and source or method of preparation. Because the value of your paper (and your reputation) can be damaged if your results are not reproducible, you must describe research materials with great care.

\subsection{Methods}
For Methods, the usual order of presentation is chronologically. Obviously, however, related methods should be described together, and straight chronological order cannot always be followed.  Be precise. Methods are similar to cookbook recipes. If a reaction mixture was heated, give the temperature. Questions such as "how" and "how much" should be precisely answered by the author and not left for the reader to puzzle over. If your method is new (unpublished) you must provide all of the needed detail. If, however, the method has been published in a journal, the literature reference is needed.

\subsubsection{Numbers, Equations and Statistics}
Spell out any number that begins a sentence, title, or heading - or better still, reword it to place the number later in the text. In general, use Arabic numerals (10, 11, 12) when referring to whole numbers 10 and above, and spell out whole numbers below 10. There are some exceptions to this rule:
\begin{itemize}
	\item If small numbers are grouped with large numbers in a comparison, use numerals, e.g.~7, 8, 10, and 13 trials; but not when numbers are used for different purposes, e.g.~10 items on each of four surveys.
	\item Numbers in a measurement with units, e.g.~6 cm, 5-mg dose, 2\%.
	\item Numbers that represent time, dates, ages, sample or population size, scores, or exact sums of money.
	\item Numbers that represent a specific item in a numbered series, e.g.~Fig.~5.
\end{itemize}

\noindent Space and punctuate equations and formulas as you would with regular words. Example: in the equation $3x + y = z$, we see that etc. Put simple equations in the text. Set off complicated equations from the text. Of course LaTeX will do so for you, like in the example here,
\begin{equation}
\x = \int_0^\infty \sqrt{1+x}\cdot\frac{\log x}{2\pi}\cdot dx.
\end{equation}
Nice, no?

Do not make the common error of mixing some of the Results in this section. There is only one rule for a properly written Materials and Materials section: Enough information must be given so that the experiments could be repeated by a competent colleague.

\section{Results}
\label{sec:results}

So now we come to the core of the paper, the data. This part of the paper is called the Results section.

There are usually two ingredients of the Results section. First, you should give some kind of overall description of the experiments, providing the "big picture" without, however, repeating the experimental details previously provided in Materials and Methods. Second, you should present the data . Most importantly, in the manuscript you should present representative data rather than endlessly repetitive data.
 
As an example, Fig.~\ref{fig:res} is included to show you how a picture can be placed in the text. Do not forget that every figure deserves a caption, placed directly below the figure. You should refer to that figure in your text as well.
%
\begin{figure}[htb]
  \centering
\centerline{\epsfig{figure=watertemp.eps,width=8cm}}
\caption{Example of an encapsulated postscript figure.}
\label{fig:res}
\end{figure}
%

Also .jpg-files can be included, as can be seen from Fig.~\ref{chaos}.
\begin{figure}[!ht]
\begin{center}
\includegraphics[width=0.75\columnwidth]{chaos.jpg}
\caption{Example of a .jpg picture.}
\label{chaos}
\end{center}
\end{figure}
%

Tables have their caption on top, as shown in table \ref{tabex}.
%
\begin{table}[ht]
\caption{Nonlinear Model Results} % title of Table
\label{tabex} % is used to refer this table in the text
\centering % used for centering table
\begin{tabular}{c c c c} % centered columns (4 columns)
\hline\hline %inserts double horizontal lines
Case & Method\#1 & Method\#2 & Method\#3 \\ [0.5ex] % inserts table
%heading
\hline % inserts single horizontal line
1 & 50 & 837 & 970 \\ % inserting body of the table
2 & 47 & 877 & 230 \\
3 & 31 & 25 & 415 \\
4 & 35 & 144 & 2356 \\
5 & 45 & 300 & 556 \\ [1ex] % [1ex] adds vertical space
\hline %inserts single line
\end{tabular}
\end{table}
%

Back to the Results. Your results should be presented in the past tense.  Although the Results section is the most important part, it is often the shortest, particularly if it is preceded by a well-written Materials and Methods section and followed by a well written Discussion.
 
The Results need to be clearly and simply stated because it is the Results that constitute the new knowledge that you are contributing to the world. The earlier parts of the paper (Introduction, Materials and Methods) are designed to tell why and how you got the Results; the later part of the paper (Discussion) is designed to tell what they mean. Obviously, therefore, the whole paper must stand or fall on the basis of the Results. Thus, the Results must be presented with crystal clarity. 

\section{Discussion}
\label{sec:discussion}
The Discussion is harder to define than the other sections. Thus, it is usually the hardest section to write. Many, if not most, Discussion sections are too long and verbose.
 
What are the essential elements of a good Discussion? The main components will be provided if the following injunctions are heeded.
\begin{enumerate}
	\item	Show relationship among observed facts. This is the primary purpose of the Discussion.
	\item Try to present the principles, relationships, and generalizations shown by the Results. And bear in mind, in a good Discussion, you discuss -you do not recapitulate - the Results.
	\item Point out any exceptions or any lack of correlation and define unsettled points. Never take the high-risk alternative of trying to cover up or fudge data that do not quite fit.
	\item Show how your results and interpretations agree (or contrast) with previously published work.
	\item Don't be shy; discuss the theoretical implications of your work, as well as possible practical applications.
	\item State your conclusions as clearly as possible.
	\item Summarize your evidence for each conclusion.
\end{enumerate} 

\noindent A discussion is also a place for reflection and for advice/warnings to other researchers. Here you can write how the research could still be improved and what future researchers should focus on. If this is a particularly long section, it could also become a new part called 'Future work'.

Much as the Methods and the Results should correspond to each other, the Introduction and the Discussion should function as a pair. Be sure this part answers what the Introduction asked. 

\section{Conclusions}
\label{sec:conclusions}
A separate "Conclusions" section should  be incorporated after the discussion. A conclusion is again  a summary of your work, a reader that only reads the conclusion should still get the main contribution of your work. A good model of a conclusion is to think of it as an extended abstract in which you can put more stress on the results and why your work is important. You want to end your paper with highlighting the significance of your work.

\section{Acknowledgements}
\label{sec:acknow}
The main text of a scientific paper is usually followed by two additional sections, namely, the Acknowledgements and the References.
As to Acknowledgements, the important element is simple courtesy. There isn't anything really scientific about this section of the paper. Two possible ingredients require consideration. First, you should acknowledge any significant technical help that you received from any individual, whether in your laboratory or elsewhere. You should also acknowledge the source of special equipment, cultures or other materials. Second, it is usually the Acknowledgements wherein you should acknowledge any outside financial assistance, such as grants, contracts or fellowships.

% To start a new column (but not a new page) and help balance the last-page
% column length use \vfill\pagebreak.
% -------------------------------------------------------------------------
%\vfill
%\pagebreak

\vspace{2em}
\noindent 
\textbf{Note on the references (remove this!)}. The reference section is the last section of a scientific paper. 
Papers are not referred to by footnotes but are cited within the body of the text (see below) \cite{day:book06}.
This way it is clear what information is obtained from what source(s) \cite{allay:book95}.
Notice that we use the APA standard for the references, meaning that the apacite package is used. When using numbered citations, other packages might be better suited.

%\renewcommand{\refname}{REFERENCES}
\bibliographystyle{apacite}
\bibliography{mybib}



\end{document}
